\documentclass[11pt]{article}

\usepackage[T1]{fontenc}
\usepackage{lmodern}
\usepackage{amsmath,amssymb,amsthm}
\usepackage[margin=1in]{geometry}
\usepackage{enumitem}
\usepackage{xcolor}
\usepackage{microtype}
\usepackage[hidelinks]{hyperref}

\newcommand{\R}{\mathbb{R}}
\newcommand{\N}{\mathbb{N}}
\newcommand{\Om}{\Omega}
\newcommand{\norm}[1]{\left\lVert #1\right\rVert}
\newcommand{\abs}[1]{\left\lvert #1\right\rvert}
\DeclareMathOperator{\dist}{dist}

\title{Paper 3, Theorem 2.2: Statement Amendment}
\author{Internal research note for the \texttt{ShenWork} formalization}
\date{12 July 2026}

\begin{document}
\maketitle

\begin{abstract}
The nonlinear stability clause printed in Theorem~2.2 of Chen--Ruau--Shen,
Part~II, asserts a $C^{1}$ exponential estimate valid for \emph{all} $t\ge 0$
from initial data small only in $L^{\infty}$.  Under the theorem's own
initial-data hypothesis this clause is not well posed at $t=0$: the $C^{1}$
norm need not be finite, and an $L^{\infty}$ ball contains functions of
arbitrarily large $C^{1}$ norm.  The argument in fact proves two correct
statements---an all-time estimate in the fractional-power (strong) norm and an
\emph{eventual} $C^{1}$ estimate for $L^{\infty}$-small data.  This note records
the explicit counterexample to the printed clause, locates the gap, gives the
corrected two-stage formulation, and states the corresponding Lean obstructions.
This is the second statement-level issue found while formalizing the series; the
first is the parameter error in Theorem~1.2 of Part~I (companion note).
\end{abstract}

\section{Source checked}

The source examined is the original submission, not a repository paraphrase:
\begin{quote}
Le Chen, Ian Ruau, and Wenxian Shen, \emph{Chemotaxis models with
signal-dependent sensitivity and a logistic-type source, II: Persistence and
stabilization}, arXiv:2604.02599v1, 3~April 2026.
\end{quote}
Theorem~2.2 there is titled ``Linear stability and instability.''  For the
positive constant equilibrium $(u_{*},v_{*})$ with $u_{*}=(a/b)^{1/\alpha}$ and
$a,b>0$, part~(1) states:
\begin{enumerate}[label=(\roman*)]
\item if $\chi_{0}<\chi^{*}_{a,b,\beta}(u_{*})$ the equilibrium is linearly
  stable;
\item if $\chi_{0}>\chi^{*}$ it is linearly unstable;
\item moreover, in the stable regime there exist $C,\delta,\lambda>0$ such that
  every $u_{0}$ satisfying the initial-data condition (1.8) \emph{and} a
  smallness condition obeys the exponential estimate (2.12) \emph{for all
  $t\ge 0$} in the $C^{1}$ norm.
\end{enumerate}
The initial-data hypothesis (1.8) requires only $u_{0}\in C(\overline{\Om})$,
positive.  Proposition~1.1 provides convergence to $u_{0}$ in $L^{\infty}$ as
$t\to 0^{+}$ and classical spatial regularity only for \emph{positive} time.
There is no $C^{1}$, Sobolev, or fractional-power smallness hypothesis on
$u_{0}$ in the statement of Theorem~2.2.

\section{The over-statement in the printed (2.12)}

The printed nonlinear clause asserts, for data small only in $L^{\infty}$, a
$C^{1}$ exponential estimate holding at \emph{every} $t\ge 0$, including $t=0$:
\[
  \norm{u(t)-u_{*}}_{C^{1}(\overline{\Om})}
    \le C\,e^{-\delta t}\,\norm{u_{0}-u_{*}}_{\ast},
  \qquad t\ge 0 .
\]
This is not well posed under the theorem's own hypothesis:
\begin{itemize}[nosep]
\item $u_{0}$ need only lie in $C(\overline{\Om})$, so $\norm{u_{0}-u_{*}}_{C^{1}}$
  may be infinite and the $t=0$ instance meaningless;
\item even for $C^{1}$ data, an $L^{\infty}$ ball contains functions of
  arbitrarily large $C^{1}$ norm, so no single $C$ controls the $t=0$
  (or $t\to 0^{+}$) $C^{1}$ size uniformly over the admitted class.
\end{itemize}

\paragraph{Explicit counterexample.}
On the unit interval $\Om=(0,1)$ with Neumann conditions, take
\[
  u_{0,N}(x) = u_{*} + N^{-1/2}\cos(N\pi x), \qquad N\in\N .
\]
For large $N$ this is positive and satisfies (1.8), with
\[
  \norm{u_{0,N}-u_{*}}_{L^{\infty}} = N^{-1/2}\to 0,
  \qquad
  \norm{u_{0,N}-u_{*}}_{C^{1}} = N^{1/2}\pi\to\infty
\]
(and, in the minimal model, exact mass $u_{*}$).  Hence no uniform-in-data
$t=0$ estimate of the form (2.12) can hold: the smallness is measured in
$L^{\infty}$, the conclusion in $C^{1}$, and the two are unrelated at $t=0$.

\section{What the published proof actually establishes}

Section~5.1 proves the result in two distinct stages, both correct.
\begin{description}
\item[Stage A (strong norm, all time).]  For data small in the fractional-power
  space $X^{\alpha}=D\big((I-L)^{\alpha}\big)$, which in one dimension embeds
  into $C^{1}$ for $\alpha>3/4$, the standard sectorial nonlinear-stability
  theorem (Henry, \emph{Geometric Theory of Semilinear Parabolic Equations},
  Ch.~5--6) gives exponential decay in $X^{\alpha}$ for all $t\ge 0$.  This is
  instantaneous and correct.
\item[Stage B ($L^{\infty}$-small, eventual).]  For data small only in
  $L^{\infty}$, the proof first evolves to a \emph{positive} time $t_{0}$, at
  which analytic-semigroup smoothing places the solution in a small
  $X^{\alpha}$ neighborhood, and then applies Stage~A.  This yields exponential
  decay for $t\ge t_{0}$---an \emph{eventual} estimate, not a $t=0$ estimate.
\end{description}
The mechanism proves an all-time estimate in the strong norm and an eventual
estimate in $C^{1}$ for $L^{\infty}$-small data.  It does not prove the printed
(2.12) at $t=0$ from $L^{\infty}$ data; that conjunction is the over-statement.

\section{Corrected statement}

Split the printed clause into the two statements the proof supports.
\begin{enumerate}[label=(\arabic*)]
\item \textbf{Strong norm, all time.}  There exist $\alpha\in(3/4,1)$ and
  $\rho,C,\delta>0$ such that
  \[
    \norm{u_{0}-u_{*}}_{X^{\alpha}}\le\rho
    \ \Longrightarrow\
    \norm{u(t)-u_{*}}_{X^{\alpha}}\le C\,e^{-\delta t}\norm{u_{0}-u_{*}}_{X^{\alpha}}
    \quad\text{for all }t\ge 0 .
  \]
\item \textbf{$L^{\infty}$-small, eventual.}  There exist $\varepsilon,C',\delta'>0$
  and, for each admissible datum, a time $t_{0}=t_{0}(u_{0})>0$ such that
  \[
    \norm{u_{0}-u_{*}}_{L^{\infty}}\le\varepsilon
    \ \Longrightarrow\
    \norm{u(t)-u_{*}}_{C^{1}}\le C'\,e^{-\delta'(t-t_{0})}
    \quad\text{for all }t\ge t_{0} .
  \]
\end{enumerate}
The linear stability/instability dichotomy of Theorem~2.2 is unaffected and
correct; only the nonlinear clause (2.12) needs the $t=0$-versus-eventual and
$L^{\infty}$-versus-$X^{\alpha}$ distinctions made explicit.

\section{A companion parameter-fidelity point}

The sufficient stability bound $\kappa<(\sqrt{\mu}+\sqrt{a\alpha})^{2}$, with
$\kappa=\chi_{0}\gamma\nu\,u_{*}^{\gamma}(1+v_{*})^{-\beta}$, is the
\emph{continuous} minimum of $(\lambda+\mu)(\lambda+a\alpha)/\lambda$ over
$\lambda>0$.  On a bounded domain the Neumann spectrum is discrete
($\lambda_{n}=n^{2}\pi^{2}$ on $(0,1)$), so the exact stable half of
Theorem~2.2 is governed by the discrete infimum
$\min_{n\ge 1}\sigma_{n}<0$, which the continuous bound only implies and can be
strictly stronger than needed.  A faithful formalization should clear the
discrete mode condition, not stop at the continuous sufficient bound.

\section{Lean amendment}

The repository records two formal obstructions to the literal clause, both
\texttt{sorry}-free.
\begin{itemize}
\item \texttt{Paper3.not\_SectorialLocalExponentialRaw\_constant\_c1Distance}
  (\path{ShenWork/Paper3/Statements.lean}) proves the raw sectorial
  local-exponential branch false when the $C^{1}$ distance is unrelated to the
  dynamics: a constantly-$1$ $C^{1}$ distance would force $2\le C e^{-\text{rate}\,t}$
  for all $t$, impossible as the right side tends to $0$.
\item \path{ShenWork/Paper3/IntervalDomainSectorialCorrectedObstruction.lean}
  records a \emph{zero-time} obstruction: \texttt{InitialTrace} controls only the
  $t\to 0^{+}$ deleted limit, not the stored slice $u(0)$, so a global solution
  may be re-anchored to an arbitrarily large constant slice at $t=0$ without
  changing global classical solvability or the initial trace---hence no all-time
  $t=0$ $C^{1}$ bound can hold.
\end{itemize}
The corrected target used in the repository is the eventual, equilibrium-specific,
full-mode orbit bound
\texttt{IntervalDomainSpectralSemigroupOrbitBoundEventualEquilibriumWithoutMass}
(linearized multiplier $e^{-d_{k}t}$ with
$d_{k}=a\alpha+\lambda_{k}-\kappa\lambda_{k}/(\lambda_{k}+\mu)$ for every mode
$k$ including $k=0$, and a uniform positive delay $t_{0}$), together with
\texttt{LocallyExponentiallyStableFromSup} for the positive equilibrium (the mass
constraint is needed only in the minimal $a=b=0$ model).  The over-stated
\texttt{...Raw} frontier---the pure Neumann heat multiplier quantified over all
$t\ge 0$---is not the correct analytic object and is refuted for a genuine
$C^{1}$ distance from $L^{\infty}$ data.

\section{Relation to Part~I (arXiv:2512.14858) Theorem 1.2}

This is the second statement-level issue found while formalizing the
Chen--Ruau--Shen series.  The first (companion note) is a genuine
\emph{parameter} error: Theorem~1.2 of Part~I omits the guard excluding
$a>0,\ b=0$, a branch in which the solution is provably unbounded
($M'(t)=aM(t)$).  The present issue is a \emph{statement over-statement}: the
intended stability result of Part~II holds in the corrected two-stage form
above, but the printed (2.12)---a $t=0$ $C^{1}$ estimate from $L^{\infty}$-small
data---is not well posed as written.

\end{document}
